% Appendix D -- How Turbo Initech is written.  The source techniques of the
% compiler itself, each shown in a small runnable program written in the
% same manner (not a copy of the compiler's source).  Sources (Law 1):
% os/tps/tps.pas (PoolChars/PoolAdd name pool, PeekChar sentinel accessor,
% insertion-ordered symbol arrays with local-first lookup, the two-chunk
% rolling source window of ChunkMax bytes, the fail-loud capacity checks);
% ADR-0007 DEC-02 (static arenas, fixed capacity, fail loud), DEC-03 (the
% three order-independence idioms), DEC-04 (insertion-ordered tables).
% Programs replayed by test-pubs-turbo.
\Gates{test-tps-lex-fpc, test-tps-type-fpc, test-tps-gen-fpc, test-compiler, test-pubs-turbo}
\chapter{How Turbo Initech is written}
\ix{compiler, techniques}

Turbo Initech is written in the language this book defines, and its own
source is the largest Turbo Initech program there is.  This appendix shows
the techniques it is built from, each in a small program written in the
same manner.  They are useful in any program that must manage tables
without pointers.

\section{A name pool}
\ix{name pool}
The compiler stores the spelling of each different identifier once, in one
large character array, and refers to a name by its position.  Comparing two
names becomes comparing two numbers once each name has been entered:

\Program{pool}

\ProgramOutput{pool}

The second \code{'total'} is found, not added again.  The tables have fixed
sizes, and running out of room sets a flag instead of overwriting memory;
the compiler reports the same condition as \code{char pool exceeds
POOLMAX=} or \code{name table exceeds NAMEMAX=} (Appendix A).

\section{A sentinel accessor}
\ix{sentinel}
Because both operands of \rw{and} are evaluated, a scanner cannot write
\code{(At <= Length(Text)) and (Text[At] = C)}.  The compiler's lexer reads
its characters through a function that returns character 0 beyond the end,
so the loops that use it need no test of the bounds:

\Program{peek}

\ProgramOutput{peek}

\section{An insertion-ordered table}
\ix{symbol table}
The compiler's symbol table is a set of parallel arrays to which rows are
only ever appended, in the order of the declarations.  In the program
below a lookup searches from the newest row backward, so a local
declaration, entered after the globals, hides a global one, and the rows of
a routine's locals are dropped at its end by moving the end of the table
back.  The compiler itself keeps every row, each marked with the routine it
belongs to, and looks a name up first among the rows of the current
routine, then among the globals; the effect is the same, and the symbol
listing (Chapter 10) can show every scope:

\Program{symtab}

\ProgramOutput{symtab}

Because rows are kept in declaration order and never hashed, the symbol
listing (Chapter 10) and the order of the generated code are the same every
time.

\section{Reading a file in blocks}
\ix{reading in blocks}
The compiler reads its source \code{ChunkMax} bytes at a time, 255 in the
compiler, keeping only two blocks in memory.  The same loop counts the lines
of a file:

\Program{chunks}

\ProgramOutput{chunks}

The loop ends when \std{blockread}{BlockRead} returns fewer bytes than it
was asked for.
